\documentclass{sfs-2487-2024}

% Monisivuinen raportti: sisällysluettelo heti pääotsikon jälkeen (6.10),
% taulukko kuvatekstillä taulukon yläpuolella (6.5.1), alaviite (6.9),
% kolme otsikkotasoa (6.4) ja perusmetatiedot jatkosivun ylätunnisteessa.

\logo{\includegraphics{logo-organisaatio}}
\doctype{Raportti}
\date{31.5.2024}
\author{Virve Virtanen}
\docid{Dnro 45/2024}

\subject{Digiprojekti}
\title{Asiakaspalautekyselyn tulokset 2024}

\contactinfo{Organisaatio Oy}{%
  Katuosoite\\
  Postinumero Postitoimipaikka\\
  Puhelinnumero\\
  www-osoite}

\begin{document}

\maketitle

\tableofcontents

\section{Johdanto}

Tässä raportissa esitetään kevään 2024 asiakaspalautekyselyn tulokset
ja niiden perusteella ehdotettavat toimenpiteet. Kysely toteutettiin
maalis--huhtikuussa 2024 verkkokyselynä\footnote{Kysely oli avoinna
18.3.--12.4.2024, ja siihen vastasi 412 asiakasta.}, ja se lähetettiin
kaikille uutiskirjeen tilaajille.

\section{Kyselyn tulokset}

\subsection{Vastaajamäärät kanavittain}

Palautetta kerättiin kolmessa kanavassa. Suurin osa vastauksista
saatiin verkkolomakkeella, mutta myös puhelinpalvelun ja
asiakastapaamisten kautta kirjattiin palautetta.

\captionof{table}{Vastaukset palautekanavittain}
\begin{tabular}{@{}lrr@{}}
  \textbf{Kanava} & \textbf{Vastauksia} & \textbf{Osuus} \\
  Verkkolomake    & 318                 & 77~\%          \\
  Puhelinpalvelu  & 64                  & 16~\%          \\
  Asiakastapaaminen & 30                & 7~\%           \\
\end{tabular}

\subsection{Keskeiset havainnot}

\subsubsection{Verkkosivuston käytettävyys}

Selvästi yleisin kehitystoive koski verkkosivuston etusivua:
yhteystietoja ja ajankohtaisia tiedotteita pidettiin vaikeasti
löydettävinä. Tuotehinnaston päivitystiheys sai kritiikkiä, ja
vastaajista valtaosa toivoi hinnaston päivittämistä vähintään
kahdesti kuukaudessa.

\subsubsection{Asiakaspalvelun tavoitettavuus}

Puhelinpalvelun jonotusaikoihin oltiin pääosin tyytyväisiä.
Sähköpostitse lähetettyjen yhteydenottojen vastausaikaa pidettiin
kuitenkin liian pitkänä, ja vastauksen saamiseen kului useissa
tapauksissa yli viikko.

\section{Toimenpide-ehdotukset}

Tulosten perusteella ehdotetaan, että etusivu uudistetaan siten,
että yhteystiedot näkyvät heti etusivulla, tuotehinnasto päivitetään
kaksi kertaa kuukaudessa ja sähköpostiyhteydenotoille asetetaan
kolmen arkipäivän vastausaikatavoite.

Ehdotusten toteutus aikataulutetaan digiprojektin ohjausryhmässä
kesäkuussa 2024. Uudistusten vaikutusta seurataan syksyllä 2024
toistettavalla kyselyllä, jonka tulokset raportoidaan johtoryhmälle
vuoden loppuun mennessä.

\vspace{\parskip}

\attachments{Kyselylomake\\Vastausten jakaumat kysymyksittäin}

\distribution{Digiprojektin ohjausryhmä}

\forinformation{Johtoryhmä}

\vspace{\parskip}

\makecontactinfo

\end{document}
